\begin{proof}
Use a state for every positive-mass prefix of length at most $L$, initialized at
$\varepsilon$. At state $u$, retain the source's stopping and letter
probabilities. For $|u|<L$, route letter $a$ to $ua$. For
$u=u_1\cdots u_L$, route it to $u_2\cdots u_La$. The latter source
state and the true successor are obtained by the \emph{same} sequence
of $L$ normalized letter maps, applied to $x_0$ and $x_{u_1}$.
The first map bounds their distance by $\Delta$, and the remaining
$L-1$ maps contract it to $D=\Delta\gamma^{L-1}$.

To pass to probability laws, take cone-order bounds
$m y\le_{\mathcal C}x\le_{\mathcal C}M y$. Positivity implies
$m f_y\le f_x\le M f_y$. Normalization gives $m\le1\le M$;
from $M/m\le e^D$ one gets $M\le e^D$ and hence
$\|f_x-f_y\|_1\le2(e^D-1)$, by summing the positive part.
Limits give the same bound when order extrema are not attained.
No coordinate $\ell^1$ bound on an arbitrary cone is used.

Each row defect is at most $2(1-s)(e^D-1)$. The constructed automaton
stops with probability at least $s$ at every visit, so its occupation
sum is at most $1/s$. The occupation-flow certificate gives
\eqref{II:U:eq:trailing-tv}. For pruning, first merge duplicate residuals,
retain the distinct vertices, and express the initial residual and each
routed successor as convex mixtures of retained vertices. The same row
defect and stopping estimates apply. Keeping the empty-prefix state
instead is always permitted within the original $N_d(L)$ budget.
\end{proof}
